\chapter{ESPECIFICACIONES TECNICAS DE MONTAJE}

\section{Alcance}

Este capitulo describe el procedimiento general para ejecutar la instalacion electrica interior de la vivienda, de acuerdo con los planos, metrados y especificaciones.

\section{Replanteo}

Antes del montaje se debe:

\begin{itemize}
  \item Revisar el plano de ambientes.
  \item Marcar ubicacion de tablero, interruptores, tomacorrientes y luminarias.
  \item Verificar recorridos de tuberias y cajas.
  \item Coordinar pases por muros, techos o losas.
\end{itemize}

\section{Canalizaciones y cajas}

\begin{itemize}
  \item Instalar tuberias segun el plano de circuitos.
  \item Evitar curvas excesivas que dificulten el cableado.
  \item Fijar cajas en posiciones firmes y niveladas.
  \item Separar circuitos cuando el diseno lo requiera.
\end{itemize}

\section{Cableado}

\begin{itemize}
  \item Cablear fase, neutro y tierra segun circuito.
  \item Evitar empalmes dentro de tuberias.
  \item Realizar empalmes solo en cajas accesibles.
  \item Identificar circuitos en tablero y cajas cuando sea posible.
\end{itemize}

\section{Tablero}

\begin{itemize}
  \item Montar el tablero en zona accesible y segura.
  \item Instalar interruptor principal, interruptor diferencial y protecciones de circuitos.
  \item Separar barra de neutro y barra de tierra.
  \item Rotular cada circuito.
\end{itemize}

\section{Puesta a tierra}

\begin{itemize}
  \item Instalar electrodo en ubicacion accesible para inspeccion.
  \item Colocar caja de registro.
  \item Conectar conductor de tierra al tablero.
  \item Verificar continuidad de conductor de proteccion.
\end{itemize}

\section{Pruebas}

Antes de energizar:

\begin{itemize}
  \item Continuidad de conductores.
  \item Aislamiento basico, si se cuenta con instrumento.
  \item Polaridad de tomacorrientes.
  \item Funcionamiento de interruptores.
  \item Identificacion de circuitos.
  \item Prueba de interruptor diferencial, si se instala.
\end{itemize}

\section{Seguridad}

\begin{itemize}
  \item Trabajar sin tension.
  \item Usar herramientas aisladas.
  \item No energizar circuitos incompletos.
  \item Mantener tablero cerrado y rotulado.
\end{itemize}
